\begin{summarybox}
\textbf{椭圆外一点到椭圆的距离极值——快记}

\begin{center}
\textbf{“椭圆先设点，距离平方算；求导找候选，比较定极值。”}
\end{center}

\begin{itemize}
  \item \textbf{设点：}
  $P=(a\cos\theta,b\sin\theta)$。
  \item \textbf{距离平方：}
  $F=(a\cos\theta-x_0)^2+(b\sin\theta-y_0)^2$。
  \item \textbf{驻点方程：}
  $(a^2-b^2)\sin\theta\cos\theta
   -ax_0\sin\theta+by_0\cos\theta=0$。
  \item \textbf{求解：}轴上外点优先化二次式；
  一般外点可用 $t=\tan(\theta/2)$ 化为四次以内方程。
  \item \textbf{收尾：}查遍实根对应点与 $(-a,0)$，
  比较距离平方，再开方。
\end{itemize}

\textbf{核心区别：}圆能直接“中心距加减半径”；
椭圆一般必须求距离函数的最值。

\medskip\textbf{法向量快记：}写切线 $xX/a^2+yY/b^2=1$，取一般式系数，得 $\vec n=(x/a^2,y/b^2)$；坐标轴顶点也适用。
\medskip\textbf{几何快记：距离取极值，连线找法线。}
\[
\overrightarrow{PM}=\lambda\left(\frac{x}{a^2},\frac{y}{b^2}\right).
\]
最近点：列 $G(\lambda)=1$，取唯一正根，还原坐标再求距离。
最远点：找全部有效法线候选点并比较；若 $\lambda<-a^2$，可直接判定唯一最远点。
轴上外点优先考虑二次函数法，一般外点优先理解法线法，万能代换保留为备用。

\end{summarybox}
